% Chương 7: Kết luận
\chapter{KẾT LUẬN}

\section{Tổng kết quá trình thiết kế}

Báo cáo đã tiếp cận bài toán nền tảng FilmLab theo quy trình thiết kế cơ sở dữ liệu từ yêu cầu nghiệp vụ đến mô hình khái niệm, mô hình logic và định hướng vật lý. Sáu Core Flow được sử dụng làm căn cứ nhận diện dữ liệu cần quản lý, gồm tìm kiếm Film Lab, đặt dịch vụ, theo dõi quá trình xử lý phim, giao file và lưu trữ ảnh số, giao dịch thiết bị, chia sẻ tri thức và hỗ trợ AI.

Trên cơ sở đó, các thực thể được chia thành những nhóm: người dùng và phân quyền; Film Lab và dịch vụ; đơn hàng và xử lý phim; kho ảnh số; chợ giao dịch; tri thức, cộng đồng và AI. Cách phân chia theo module giúp từng thành viên phát triển phần phụ trách, đồng thời vẫn duy trì liên kết xuyên suốt thông qua các thực thể trung tâm như \entity{users}, \entity{film_labs} và \entity{orders}.

\section{Kết quả đạt được}

\begin{itemize}
    \item Xác định tác nhân, yêu cầu dữ liệu, quy tắc nghiệp vụ và các mối quan hệ quan trọng của hệ thống.
    \item Đồng bộ phạm vi P1--P5 ở Chương 2--3, lập danh mục thực thể, bảng cardinality và sơ đồ truy vết Order--Processing--Archive theo mô hình hiện tại. ERD tổng thể cuối cùng vẫn cần hoàn thiện sau khi bổ sung các quan hệ còn thiếu.
    \item Đồng bộ 36 bảng hiện có giữa mô hình logic và DDL: 11 bảng Core, 5 bảng P3, 14 bảng P4 và 6 bảng P5. Bảng dùng chung chỉ định nghĩa một lần trong từng chương; phân tích khóa ứng viên và các giới hạn chuẩn hóa còn tồn tại.
    \item Đề xuất định hướng thiết kế vật lý trên PostgreSQL về kiểu dữ liệu, ràng buộc, chỉ mục, audit, soft-delete, multi-tenant và lưu trữ file scan.
\end{itemize}

\section{Đánh giá thiết kế}

Phần Core Domain P1 đã bổ sung phụ thuộc hàm và khóa ứng viên cho các
quan hệ lõi, ràng buộc chọn đúng một dịch vụ hoặc một gói trên mỗi hạng
mục và audit chung cho cả bảng nối. Subtotal và total amount được tính
từ giá, số lượng, phí và giảm giá đã chốt thay vì lưu lặp trong orders.
Nội dung gói đã được đặt được bảo toàn để giữ khả năng truy vết.
Các quy tắc giao dịch và kiểm tra liên bảng đã được đặc tả; việc triển
khai, kiểm thử trigger và tích hợp đầy đủ các module vẫn cần hoàn thiện.

Thiết kế bảo đảm khả năng truy vết dữ liệu từ người dùng, Film Lab và đơn hàng đến quá trình xử lý và kho ảnh số. RBAC hỗ trợ một người dùng có nhiều vai trò; phạm vi tenant được nhận diện thông qua Film Lab; lịch sử nghiệp vụ quan trọng được giữ lại thay vì xóa cứng. Giá dịch vụ tại thời điểm đặt được lưu như một snapshot, giúp đơn hàng không bị thay đổi khi danh mục giá được cập nhật.

Việc tách file ảnh khỏi cơ sở dữ liệu và lưu URI cùng metadata phù hợp với đặc điểm dữ liệu scan có dung lượng lớn; checksum và versioning chưa có trong mô hình hiện tại. Các nhóm dữ liệu hành vi, tri thức và hội thoại tạo nền tảng để mở rộng recommendation, semantic search và trợ lý RAG trong tương lai.

\section{Hạn chế}

Báo cáo hiện tập trung vào thiết kế và DDL, chưa kiểm thử schema trên PostgreSQL, dữ liệu mẫu hoặc hiệu năng thực tế. Cần bổ sung phạm vi P2, ánh xạ hạng mục vào cuộn phim, versioning scan và các quan hệ Logistics, Payment, Notification, nguồn tri thức/AI. Chuỗi tags và bộ đếm view\_count của articles còn cần xử lý trước khi kết luận chuẩn hóa toàn phần. Trigger liên bảng, quyền truy cập và các miền trạng thái chưa chốt vẫn cần triển khai, kiểm chứng; việc 36 bảng có định nghĩa nhất quán chưa có nghĩa toàn bộ yêu cầu đã hoàn thành.

\section{Hướng phát triển}

Trong giai đoạn tiếp theo, nhóm cần hoàn thiện ERD tổng thể sau cross-review, đồng bộ tên bảng và khóa giữa tất cả module, sau đó kiểm thử mô hình bằng một bộ dữ liệu đại diện. Khi chuyển sang triển khai, hệ thống có thể bổ sung chỉ mục không gian cho tìm Film Lab, full-text hoặc vector search cho kho tri thức, phân vùng bảng log theo thời gian, chính sách Row-Level Security cho multi-tenant và cơ chế sao lưu/khôi phục định kỳ.

Nhìn chung, mô hình hiện tại đã hình thành nền tảng dữ liệu có cấu trúc cho hệ sinh thái nhiếp ảnh phim. Việc tiếp tục hoàn thiện các module còn thiếu và kiểm chứng bằng dữ liệu thực tế sẽ giúp thiết kế đạt tính nhất quán, khả năng mở rộng và độ tin cậy cần thiết trước khi được hiện thực hóa.
